\section{Билет 46. Теплоемкость. Удельная теплоемкость. Молярная теплоемкость. Теплоемкость идеального газа при изохорическом процессе с выводом. Соотношение Майера.}

\begin{definition}
\textbf{Теплоёмкостью} $C$ тела называется физическая величина, равная количеству теплоты, необходимому для нагревания этого тела на один кельвин:
\begin{equation}
    C = \frac{\delta Q}{dT}. \label{eq:heat_capacity_def}
\end{equation}
Единица измерения --- Дж/К. Теплоёмкость зависит от массы тела, поэтому вводятся удельные и молярные величины.
\end{definition}

\begin{definition}
\textbf{Удельная теплоёмкость} $c$ --- теплоёмкость единицы массы вещества:
\begin{equation}
    c = \frac{C}{m} = \frac{1}{m} \frac{\delta Q}{dT}, \quad \left[ c \right] = \frac{\text{Дж}}{\text{кг} \cdot \text{К}}.
\end{equation}
\textbf{Молярная теплоёмкость} $C_\mu$ --- теплоёмкость одного моля вещества:
\begin{equation}
    C_\mu = \frac{C}{\nu} = \frac{1}{\nu} \frac{\delta Q}{dT}, \quad \left[ C_\mu \right] = \frac{\text{Дж}}{\text{моль} \cdot \text{К}}.
\end{equation}
Связь между ними: $C_\mu = c \cdot M$, где $M$ --- молярная масса.
\end{definition}

\begin{theorem}[Теплоёмкость при изохорном процессе]
При постоянном объёме ($V = \text{const}$) работа газа $A = 0$. Из первого начала термодинамики:
\begin{equation}
    \delta Q_V = dU.
\end{equation}
Для идеального газа внутренняя энергия зависит только от температуры: $U = \frac{i}{2} \nu R T$, где $i$ --- число степеней свободы. Тогда:
\begin{equation}
    dU = \frac{i}{2} \nu R \, dT.
\end{equation}
Подставляя в определение молярной теплоёмкости \eqref{eq:heat_capacity_def}:
\begin{equation}
    C_V = \frac{1}{\nu} \frac{\delta Q_V}{dT} = \frac{1}{\nu} \frac{dU}{dT} = \frac{i}{2} R. \label{eq:Cv_ideal}
\end{equation}
\textbf{Вывод:} Молярная теплоёмкость идеального газа при постоянном объёме зависит только от числа степеней свободы молекул и универсальной газовой постоянной.
\end{theorem}

\begin{theorem}[Соотношение Майера]
При изобарном процессе ($p = \text{const}$) первое начало термодинамики даёт:
\begin{equation}
    \delta Q_p = dU + p \, dV.
\end{equation}
Разделим на $\nu \, dT$ и учтём, что $dU = \nu C_V \, dT$:
\begin{equation}
    C_p = \frac{1}{\nu} \frac{\delta Q_p}{dT} = C_V + \frac{p}{\nu} \frac{dV}{dT}.
\end{equation}
Из уравнения состояния идеального газа $pV = \nu R T$ при $p = \text{const}$:
\begin{equation}
    p \, dV = \nu R \, dT \quad \Rightarrow \quad \frac{p}{\nu} \frac{dV}{dT} = R.
\end{equation}
Окончательно получаем \textbf{соотношение Майера}:
\begin{equation}
    C_p = C_V + R. \label{eq:mayer_relation}
\end{equation}
\end{theorem}

\begin{property}
Из соотношения Майера \eqref{eq:mayer_relation} и формулы \eqref{eq:Cv_ideal} следуют выражения для молярных теплоёмкостей идеального газа:
\begin{equation}
    C_V = \frac{i}{2} R, \quad C_p = \frac{i+2}{2} R.
\end{equation}
\textbf{Коэффициент Пуассона} (показатель адиабаты) определяется как отношение теплоёмкостей:
\begin{equation}
    \gamma = \frac{C_p}{C_V} = \frac{i+2}{i}.
\end{equation}
Значения для различных газов:
\begin{itemize}
    \item Одноатомные ($i=3$): $\gamma = 5/3 \approx 1.67$;
    \item Двухатомные ($i=5$): $\gamma = 7/5 = 1.4$;
    \item Многоатомные ($i=6$): $\gamma = 8/6 \approx 1.33$.
\end{itemize}
\end{property}

\begin{remark}
\textbf{Физический смысл разности $C_p - C_V = R$:} При нагревании газа при постоянном давлении часть подводимой теплоты расходуется на совершение работы расширения. Работа одного моля идеального газа при изобарном нагревании на 1 К равна универсальной газовой постоянной $R \approx 8.31$ Дж/(моль·К). Именно эта величина составляет разницу между теплоёмкостями при постоянном давлении и постоянном объёме.
\end{remark}
